\documentclass[11pt]{article}
\usepackage[a4paper,margin=27mm]{geometry}
\usepackage{amsmath,amssymb,amsthm,microtype,hyperref}
\setlength{\emergencystretch}{2em}
\hypersetup{pdfauthor={Ethan Walker},colorlinks=true,urlcolor=blue,citecolor=blue,linkcolor=blue}
\newtheorem{theorem}{Theorem}
\newtheorem{corollary}[theorem]{Corollary}
\newtheorem{lemma}[theorem]{Lemma}
\title{Collatz prefixes with finitely many deviations\\from rational periodic trajectories}
\author{Ethan Walker\thanks{Prepared from the author's October 2026 research notes with OpenAI ChatGPT/Codex assistance, including proof reconstruction and exposition. This preprint has not undergone independent human peer review.}}
\date{21 October 2026}
\begin{document}\maketitle
\begin{abstract}
For the accelerated odd Collatz map, a prefix with at most $R$ divisions by $4$ or more has length at most $(1+\varepsilon)\log_2(n+1)+O_{R,\varepsilon}(1)$, for each fixed $R$ and $\varepsilon>0$. We prove a common generalization: the same leading coefficient bounds reference depth for prefixes with finitely many deviations from any fixed finite collection of rational periodic trajectories. The proof uses the $p$-adic Subspace Theorem and finiteness of nondegenerate $S$-unit equations. The constants are ineffective and not uniform in the number of deviations.
\end{abstract}
\section{The simplest statement}
For a positive odd integer $n$, put
\[
a(n)=v_2(3n+1),\qquad U(n)=\frac{3n+1}{2^{a(n)}}.
\]
A step is \emph{deep} if $a(n)\ge2$; otherwise it is shallow.
\begin{corollary}\label{fixed}
For every fixed integer $R\ge0$ and every $\varepsilon>0$, there is a finite constant $K_{R,\varepsilon}$ such that a prefix
$n_0,\ldots,n_h$ with at most $R$ deep steps satisfies
\[
h\le(1+\varepsilon)\log_2(n_0+1)+K_{R,\varepsilon}.
\]
\end{corollary}
The leading coefficient $1$ cannot be decreased: starting at
$n_0=2^{h+1}-1$ gives $h$ shallow steps. For $R=0$ the bound follows directly from
$n_t+1=(3/2)^t(n_0+1)$ and divisibility by $2^{h+1}$.
The work below allows any fixed number of deep steps, whose depths can be arbitrarily large.

\section{A finite periodic dictionary}
Extend $U$ to odd rational $2$-adic integers using the same valuation rule. Fix a finite nonempty union $C$ of complete periodic trajectories of this map. Write
\[
b(c)=v_2(3c+1)\ge1,\qquad g(c)=\frac{3c+1}{2^{b(c)}}\in C.
\]
These rationals have denominators coprime to $6$: the period equation has denominator dividing $2^A-3^m$. All valuations below are well defined, and $c$ and $3c+1$ are nonzero.
A reference path is any $c_0,\ldots,c_h\in C$. Its step $t$ needs a \emph{repair} if
\[
a(n_t)\ne b(c_t)\quad\hbox{or}\quad c_{t+1}\ne g(c_t).
\]
A change of reference phase also costs a repair. Define the reference depth
$B_h=\sum_{t=0}^{h-1}b(c_t)$, and let $\beta$ be the smallest mean of $b$ over the cycles in $C$.
\begin{theorem}\label{main}
Fix $C$, $R\ge0$, and $\varepsilon>0$. Suppose a positive odd orbit prefix has a reference path with at most $R$ repairs. Assume either all members of $C$ are negative, or
\[
n_t\ge L_C:=\left\lceil2\max_{c\in C}|c|\right\rceil+2
\quad(0\le t\le h).
\]
Then constants depending only on $C,R,\varepsilon$ satisfy
\[
B_h\le(1+\varepsilon)\log_2(n_0+1)+K,
\qquad
h\le\frac{1+\varepsilon}{\beta}\log_2(n_0+1)+K'.
\]
\end{theorem}
Corollary~\ref{fixed} is $C=\{-1\}$: here $b=1$ and $g(-1)=-1$. The size condition matters for positive references: the actual fixed point $1$ could otherwise follow itself forever.

\section{Two standard Diophantine inputs}
\noindent\textbf{Subspace Theorem.}
Fix a finite set of places $S$ of $\mathbb Q$ containing infinity. At each place choose $d$ independent rational linear forms in $d$ variables. For every $\eta>0$, the integer vectors $X\ne0$ satisfying
\[
\prod_{v\in S}\prod_{j=1}^d|L_{j,v}(X)|_v
\le \|X\|_\infty^{-\eta}
\]
lie in finitely many proper rational linear subspaces. This is the rational case of the $p$-adic theorem in Bilu~\cite[Theorem 2.3]{bilu}.

\medskip\noindent\textbf{$S$-unit finiteness.}
A minimal zero sum of a fixed number of terms, each a fixed nonzero rational coefficient times an element of a fixed finitely generated subgroup of $\mathbb Q^*$, has finitely many projective term ratios. Divide by one term and apply Evertse--Schlickewei--Schmidt~\cite[Theorem 1.1]{ess}. We use the group generated by $2$ and $3$.

\section{Proof of Theorem~\ref{main}}
Choose an odd common denominator $D$ for $C$, coprime to $3$. Put
\[
a_t=a(n_t),\quad A_t=\sum_{j<t}a_j,\quad
Z_t=\frac{2^{A_t}}{3^t}(n_t-c_t)>0.
\]
A good step leaves $Z_t$ unchanged. At a repair,
\begin{equation}\label{increment}
Z_{t+1}-Z_t=3^{-t-1}
\left((3c_t+1)2^{A_t}-c_{t+1}2^{A_{t+1}}\right).
\end{equation}
Moreover
\begin{equation}\label{ratio}
\frac{Z_{t+1}}{Z_t}
=\left(1+\frac1{3n_t}\right)
\frac{1-c_{t+1}/n_{t+1}}{1-c_t/n_t}.
\end{equation}
The domain assumptions bound this ratio above and away from zero uniformly. With at most $R$ repairs, $Z_t\asymp_{C,R} n_0+1$.
If all pre- and post-repair values tend to infinity along a family, all snapshot ratios $Z_t/Z_0$ tend to $1$.

\subsection{Reductions}
Prove the depth bound by induction on $R$. Without repairs,
$D(n_h-c_h)=3^h2^{-B_h}D(n_0-c_0)$.
The endpoint difference is an even $2$-adic integer, so
$v_2(D(n_0-c_0))\ge B_h+1$.
Its nonzero integral numerator has absolute value $O_C(n_0+1)$, proving the base case.

From any fixed positive starting value there is a finite maximum length of an admissible prefix with at most $R$ repairs. Indeed a good run starting at $n$ with reference $c$ has length at most $v_2(D(n-c))-1$: each good step decreases that valuation by $b(c)\ge1$. There are finitely many reference choices and possible first-repair locations. Induction over the remaining repairs proves the assertion, uniformly on finite sets of starting values.

If the bound fails for fixed $\varepsilon>0$, choose a family with
\begin{equation}\label{failure}
B_h-(1+\varepsilon)\log_2(n_0+1)\longrightarrow\infty.
\end{equation}
Replace each reference path by one needing the fewest repairs. Two paths with at most $R$ repairs have the same reference depths outside the union of their repair sets. Their $B_h$ differ by at most $2R\max b$, so failure persists.
The fixed-start assertion permits a subsequence with $n_0\to\infty$.
Every post-repair value tends to infinity as well: otherwise a post-repair value stays in a finite set; the prefix before that repair has at most $R-1$ repairs and obeys induction, while its bounded-start suffix has bounded length, contradicting~\eqref{failure}. Pre-repair values tend to infinity since $n_{t+1}\le(3n_t+1)/2$.

Pass to a subsequence with a fixed positive number of repairs, fixed reference labels at repairs, and each successive repair gap either fixed or tending to infinity. Group repairs separated by fixed gaps into clusters
$[\ell_i,k_i)$, $1\le i\le s$, with fixed widths $w_i=k_i-\ell_i$. Distinct clusters separate by gaps tending to infinity. Set $k_0=0$ and
$\delta_i=Z_{k_i}-Z_{k_{i-1}}$.
There are no changes of $Z$ between clusters.

\begin{lemma}\label{nonzero}
Every cluster increment $\delta_i$ is nonzero.
\end{lemma}
\begin{proof}
A zero increment says that the affine word of actual exponents across the cluster sends $c_{\ell_i}$ to $c_{k_i}$. Every intermediate rational has denominator dividing $D$ times a power of $2$ by forward iteration, and dividing $D$ times a power of $3$ by backward iteration. Thus its denominator divides $D$.
Backward iteration from the odd endpoint shows all intermediates are odd $2$-adic integers. Every exponent is consequently the exact valuation at its preceding intermediate. The word follows $g$ from $c_{\ell_i}$, hence stays in $C$. Replacing the internal references by this trajectory removes repairs and leaves both endpoints unchanged, contradicting minimality.
\end{proof}

\subsection{An inequality at three places}
Let $k=k_s$, $d=s+1$, and take the positive integer vector
\[
X_i=D3^kZ_{k_i}\quad(0\le i\le s),\qquad
H=\max_iX_i\ll3^h(n_0+1).
\]
At infinity use the independent forms
$X_1-X_0,\ldots,X_s-X_{s-1},X_s$; at $2$ and $3$ use the coordinate forms. Their product satisfies
\begin{equation}\label{product}
\Pi(X)\ll(n_0+1)2^{-B_h}.
\end{equation}
Here constants may depend on the fixed cluster widths. To verify it, put
$v_i=v_2(D(n_{k_i}-c_{k_i}))\ge1$. Then
\[
|X_i|_2=2^{-A_{k_i}-v_i},\qquad
|X_i|_3\le3^{-(k-k_i)}.
\]
Summing~\eqref{increment} across a cluster and using $a_t\ge1$ gives
\[
|\delta_i|\ll 2^{A_{k_i}}3^{-k_i}(3/2)^{w_i}.
\]
Indeed a summand is at most a constant times $2^{A_{t+1}}3^{-t-1}$; its ratio to $2^{A_{k_i}}3^{-k_i}$ is at most $(3/2)^{k_i-t-1}$.
Each snapshot valuation pays for the depth of the following good gap, leaving at least one unit at its endpoint. These gaps cover the complement of the clusters, so
\[
\sum_{i=0}^s v_i\ge B_h-(\max b)\sum_iw_i+d.
\]
Multiplying the real and $2$-adic and $3$-adic estimates cancels every $A_{k_i}$ and $k_i$ factor. The final real form contributes $Z_{k_s}\ll n_0+1$. This proves~\eqref{product}.

Since $h\le B_h$, failure~\eqref{failure} implies
$\Pi(X)\le H^{-\eta}$ eventually for some fixed $\eta>0$.
More explicitly, writing $L=\log_2(n_0+1)$ and $\alpha=\log_2 3$ gives
$\log_2H\le\alpha B_h+L+O(1)$ and
$\log_2\Pi\le L-B_h+O(1)$; the ratio
$(B_h-L)/(\alpha B_h+L)$ stays bounded away from zero.
The Subspace Theorem applies to these fixed forms in dimension $d$.

\subsection{Eliminating the exceptional subspaces}
One proper rational subspace contains infinitely many vectors. Thus fixed rational coefficients $z_i$, not all zero, satisfy
$\sum_{i=0}^s z_iZ_{k_i}=0$ on an infinite subfamily.
By~\eqref{ratio} and the reductions, $Z_{k_i}/Z_0\to1$; hence $\sum z_i=0$. Rewriting in increments,
\begin{equation}\label{relation}
\sum_{i=1}^s q_i\delta_i=0,\qquad q_i=\sum_{j=i}^s z_j,
\end{equation}
with some $q_i\ne0$.
Expand~\eqref{relation} by~\eqref{increment}, omitting clusters with zero $q_i$. This is a zero sum of at most $2R$ nonzero terms, each a fixed rational coefficient times $2^u3^{-t-1}$. Partition them into minimal vanishing subsums; there are finitely many partitions, so fix one on a subsequence. By $S$-unit finiteness, fix all within-subsum projective ratios on a further subsequence.
For terms at times $t,t'$, the $3$-adic valuation of their ratio is $-t+t'$ plus a fixed coefficient valuation. A fixed ratio forces $t-t'$ bounded. No minimal subsum therefore crosses between distinct clusters. Summing the subsums in a participating cluster gives $q_i\delta_i=0$, contradicting Lemma~\ref{nonzero}. The depth bound follows.

Finally a path with at most $R$ repairs consists of at most $R+1$ genuine periodic pieces and $R$ repair positions. Each piece has depth at least $\beta$ times its length minus a constant depending on $C$. Thus $B_h\ge\beta h-O_{C,R}(1)$, proving the length bound.

\section{Scope and prior literature}
This is an initial-prefix bound for each fixed repair budget. Its leading coefficient is sharp already at reference $-1$, but no effective constant is obtained. It does not control a growing repair budget, establish positive density of deep steps, or prove Collatz convergence.

Brox~\cite{brox} proves finiteness results for cycles with few descents using logarithmic forms. Those cycle results are not the finite-prefix assertion here. Sharpe's research repository~\cite{sharpe} studies necessary growth of parity shift defects; that count should not be identified with this repair count without a separate translation. Subspace methods themselves have appeared in Collatz research, for example Fleckinger and Abdoulkarim~\cite{fa}. No new Diophantine theorem or first use of these methods for Collatz is claimed.

Relative to these named results, the additional statement is the explicit finite-prefix estimate for fixed deep-step count, its fixed finite periodic-dictionary extension, and removal of exceptional subspaces. A comprehensive priority claim needs a broader specialist literature review.

\begin{thebibliography}{9}\small
\bibitem{bilu} Y. Bilu, The many faces of the Subspace Theorem (after Adamczewski, Bugeaud, Corvaja, Zannier\ldots), \emph{Ast\'erisque} 317 (2008), 1--38. \href{https://arxiv.org/abs/0907.2098}{arXiv:0907.2098}.
\bibitem{ess} J.-H. Evertse, H. P. Schlickewei and W. M. Schmidt, Linear equations in variables which lie in a multiplicative group, \emph{Ann. of Math.} 155 (2002), 807--836. \href{https://arxiv.org/abs/math/0409604}{arXiv:math/0409604}.
\bibitem{brox} T. Brox, Collatz cycles with few descents, \emph{Acta Arith.} 92 (2000), 181--188. \href{https://doi.org/10.4064/aa-92-2-181-188}{doi:10.4064/aa-92-2-181-188}.
\bibitem{sharpe} M. Sharpe, Collatz research repository, especially \texttt{paper/defects.tex}, accessed 21 October 2026. \url{https://github.com/msharpe248/collatz}.
\bibitem{fa} J.-J. Fleckinger and A. Abdoulkarim, Sur la conjecture de Collatz, preprint (2016). \href{https://arxiv.org/abs/1607.02370}{arXiv:1607.02370}.
\end{thebibliography}
\end{document}
